% =============================================================================
% Project : THE_PHILOSOPHY_OF_VIABILITY
% Book    : Book I — Ontology of Viability
% Book ID : VS.B01
% Part    : Part I — Identity and Manifestation (VS.B01.P01)
% Chapter : Chapter 07 — Formal and Logical Reconstruction
% Chapter ID : VS.B01.P01.C07
%
% File    : 4_ASSETS/diagrams/VS.B01.P01.C07.S04.SS02.P001.DG01.typed_dependency_structure.tex
%
% Section / Chapter / Asset:
%   Diagram DG01 — Typed Dependency Structure
%
% Structural Identifier : VS.B01.P01.C07.S04.SS02.P001.DG01
% Asset Type : DG
%
% Purpose :
%   To distinguish definitional, conceptual, logical, ontological, and
%   justificatory relations among the elements of a formal reconstruction.
%
% Insertion point :
%   Section 04, subsection "Dependency Graphs and the Control of Conceptual
%   Order".
%
% Representational limit :
%   The relations shown are methodologically distinct and must not be
%   collapsed into a single undifferentiated dependency relation.
%
% Dependencies :
%   TikZ; libraries used: arrows.meta, positioning. All are loaded by the
%   project preamble; no per-file package or library loading.
%
% =============================================================================

\begin{figure}[htbp]
    \centering
    \hyphenpenalty=10000 \exhyphenpenalty=10000\relax  % no hyphenation inside diagram labels

    \begin{tikzpicture}[
        font=\small,
        node distance=1.0cm and 0.7cm,
        % --- nodes ---
        concept/.style={
            draw=blue!55!black, fill=blue!8, rounded corners=3pt,
            align=center, text width=3.1cm, minimum height=1.6cm,
            inner sep=5pt, line width=0.7pt
        },
        % --- five relation types: distinct by line pattern, weight and colour ---
        rel/.style={-{Stealth[length=2.5mm]}},
        definitional/.style={rel, line width=1.8pt, draw=blue!55!black},
        conceptual/.style={rel, line width=0.9pt, draw=black!75},
        logical/.style={rel, line width=0.9pt, dashed, draw=purple!60!black},
        ontological/.style={rel, line width=1.2pt, densely dotted,
                            draw=green!40!black},
        justificatory/.style={rel, line width=0.9pt,
                              dash pattern=on 5pt off 2pt on 1pt off 2pt,
                              draw=red!55!black},
        legendtext/.style={font=\scriptsize, align=center, text width=11cm}
    ]

    % ---------------------------------------------------------------
    % Grid of concepts
    % ---------------------------------------------------------------
    \node[concept] (vocab) {Conceptual vocabulary};

    \node[concept, below=of vocab] (explication) {Conceptual\\explication};
    \node[concept, left=of explication]  (definition) {Definitional\\introduction};
    \node[concept, right=of explication] (primitive) {Primitive conceptual\\or formal roles};

    \node[concept, below=of definition]  (formal) {Formal expressions\\and definitions};
    \node[concept, below=of explication] (claims) {Substantive claims\\and propositions};
    \node[concept, below=of primitive]   (interpretation)
        {Interpretation\\and ontological\\characterization};

    \node[concept, text width=4.6cm, minimum height=1.2cm, below=of claims] (assessment)
        {Argument, assessment,\\and justificatory support};

    % ---------------------------------------------------------------
    % Typed relations
    % ---------------------------------------------------------------
    \draw[definitional] (vocab) -- (definition);
    \draw[conceptual]   (vocab) -- (explication);
    \draw[conceptual]   (vocab) -- (primitive);

    \draw[definitional] (definition)  -- (formal);
    \draw[conceptual]   (explication) -- (claims);
    \draw[conceptual]   (primitive)   -- (interpretation);

    \draw[logical]       (formal) -- (claims);
    \draw[ontological]   (interpretation) -- (claims);
    \draw[justificatory] (claims) -- (assessment);
    \draw[justificatory] (interpretation) -- (assessment);

    % ---------------------------------------------------------------
    % Legend (one line sample per relation type)
    % ---------------------------------------------------------------
    \node[below=0.8cm of assessment, legendtext] (legend) {
        \tikz[baseline=-0.5ex]\draw[definitional] (0,0) -- (0.8,0);~definitional
        \quad
        \tikz[baseline=-0.5ex]\draw[conceptual] (0,0) -- (0.8,0);~conceptual
        \quad
        \tikz[baseline=-0.5ex]\draw[logical] (0,0) -- (0.8,0);~logical\\[3pt]
        \tikz[baseline=-0.5ex]\draw[ontological] (0,0) -- (0.8,0);~ontological
        \quad
        \tikz[baseline=-0.5ex]\draw[justificatory] (0,0) -- (0.8,0);~justificatory
    };

    \node[below=0.4cm of legend, text width=11cm, align=left,
          font=\footnotesize\itshape] {
        \textbf{Interpretive constraint:}
        The arrows indicate distinct kinds of theoretical connection.
        They do not establish that every connection is necessary,
        that every path is complete, or that formal introduction
        determines ontological order. In particular, definitional
        introduction and conceptual articulation do not by themselves
        establish ontological dependence.
    };

    \end{tikzpicture}

    \caption{A typed map of definitional and conceptual dependencies. The
    relations shown are methodologically distinct and must not be
    collapsed into a single undifferentiated dependency relation.}
    \label{fig:typed-dependency-structure}
\end{figure}

% =============================================================================
% END OF DIAGRAM DG01
% =============================================================================
